\section*{第 \ref{ch:g8:nervous-communication} 章——神经通信}

\begin{solution}{exo:g8:nervous-communication:1}
\omterm{def:g8:nervous-communication:neuron}{神经元}是\omterm{def:g5:brain-in-command:system}{神经系统}的工作\omterm{def:g5:first-look-at-cells:cell}{细胞}：一个\omterm{def:g5:first-look-at-cells:cell}{细胞}体
连着一根细长的纤维，有的从\omterm{def:g4:classifying-animals:vertebrate}{脊柱}一直延伸到脚趾。\omterm{def:g8:nervous-communication:neuron}{神经}是
由数千根这样的纤维捆束而成的线缆。
\end{solution}

\begin{solution}{exo:g8:nervous-communication:2}
短促的电信号，速度最高可达\qty{100}{m/s}。
信号彼此相同，强度以其发放速率编码——
从轻柔的敲击到急促的鼓点。
\end{solution}

\begin{solution}{exo:g8:nervous-communication:3}
感觉\omterm{def:g8:nervous-communication:neuron}{神经元}：传递信息的线路，从\omterm{def:g1:the-five-senses:sense}{感觉器官}到\omterm{def:g5:brain-in-command:system}{脊髓}和脑。
运动\omterm{def:g8:nervous-communication:neuron}{神经元}：指令线路，从\omterm{def:g5:brain-in-command:system}{脊髓}到\omterm{def:g3:bones-and-muscles:muscle}{肌肉}。
中间\omterm{def:g8:nervous-communication:neuron}{神经元}：负责决策，位于脑和\omterm{def:g5:brain-in-command:system}{脊髓}之内。
\end{solution}

\begin{solution}{exo:g8:nervous-communication:4}
\omterm{def:g8:nervous-communication:synapse}{突触}是微观的间隙，一个\omterm{def:g8:nervous-communication:neuron}{神经元}的纤维末端
在此与下一个\omterm{def:g5:first-look-at-cells:cell}{细胞}相接。化学信使穿过间隙并与受体对接——
若到达的量足够，
便会在下一个\omterm{def:g8:nervous-communication:neuron}{神经元}中重新启动信号。
\end{solution}

\begin{solution}{exo:g8:nervous-communication:5}
可塑性（\omterm{def:g8:nervous-communication:synapse}{突触}会因使用而增强，
不用则减弱，即学习）；单向传递（递质仅在一侧释放）；
以及加权整合（每个\omterm{def:g8:nervous-communication:neuron}{神经元}在发放前，对数千个兴奋性
与抑制性输入加和）。
\end{solution}

\begin{solution}{exo:g8:nervous-communication:6}
\omterm{def:g8:nervous-communication:neuron}{神经}：毫秒级，精确指向单一目标，专用线路，
效应随信号消失。\omterm{def:g8:puberty:hormones}{激素}：数分钟到数年，
遍及全身，经\omterm{prop:g4:journey-of-food:absorb}{血液}广播，效应持续
只要信使在\omterm{prop:g4:journey-of-food:absorb}{血液}中循环便持续存在。
\end{solution}

\begin{solution}{exo:g8:nervous-communication:7}
眼睛的感觉\omterm{def:g8:nervous-communication:neuron}{神经元}把球的位置传送到
脑；中间\omterm{def:g8:nervous-communication:neuron}{神经元}在\omterm{def:g8:nervous-communication:synapse}{突触}处对信号流加权——
所计算的路径即逐个连接处的输入求和；
运动\omterm{def:g8:nervous-communication:neuron}{神经元}传出指令，其速率决定各\omterm{def:g3:bones-and-muscles:muscle}{肌肉}的力量，
用于迈步、前倾，
以及手部的握合。
\end{solution}

\begin{solution}{exo:g8:nervous-communication:8}
每一次练习都让信号经过同样的\omterm{def:g8:nervous-communication:synapse}{突触}，
被经过的\omterm{def:g8:nervous-communication:synapse}{突触}得到增强：一个月后，
抛接球的信号路径已预先加权，原本需要
费力加权的过程如今已自动运行。这种持久性
正是\omterm{def:g8:nervous-communication:synapse}{突触}增强的结果。
\end{solution}

\begin{solution}{exo:g8:nervous-communication:9}
因为\omterm{def:g8:nervous-communication:synapse}{突触}是信号转变为化学过程之处——
释放的递质、对接位点，间隙向
到达的任何分子开放。电线中的电流很难
被仿造；化学信使则可被模仿、阻断
或耗尽——因此麻醉剂、毒液与
药物都在此起作用。
\end{solution}

\begin{solution}{exo:g8:nervous-communication:10}
(a) 毫秒级，作用于单臂的\omterm{def:g3:bones-and-muscles:muscle}{肌肉}，经有线通路，
效应随即消失：\omterm{def:g8:nervous-communication:neuron}{神经}调节。(b) 数年，遍及全身，经\omterm{prop:g4:journey-of-food:absorb}{血液}传递，
持续数个季节：\omterm{def:g8:puberty:hormones}{激素}调节。(c) 接力过程：先由\omterm{def:g8:nervous-communication:neuron}{神经}瞬时启动，
再由\omterm{prop:g4:journey-of-food:absorb}{血液}中的化学信使接续，在惊吓之后维持警觉数分钟——
二者依次发生。
\end{solution}

\begin{solution}{exo:g8:nervous-communication:11}
指令到达连接处却无法通过：\omterm{def:g3:bones-and-muscles:muscle}{肌肉}
得不到任何信号。受害者逐渐瘫痪——
意识、感觉和决策能力仍在，但回路在其
最后一个环节（即从指令到\omterm{def:g3:bones-and-muscles:muscle}{肌肉}）被切断。
\end{solution}

\begin{solution}{exo:g8:nervous-communication:12}
学习：技能即经使用而增强的\omterm{def:g8:nervous-communication:synapse}{突触}——
抛接球者的一个月，就储存在这些连接之中。决策：每个\omterm{def:g8:nervous-communication:neuron}{神经元}
依据其各连接处的输入之和决定是否发放——
回路中的“决定”即是连接处的运算。脆弱性：
麻醉剂、毒液和药物作用于信号转化为
化学物质之处——即间隙。电线只负责传导；连接
会学习、会选择，也会被欺骗。
\end{solution}

\begin{solution}{pb:g8:nervous-communication:1}
\textbf{1.} 肌腱的拉伸激活感觉\omterm{def:g8:nervous-communication:neuron}{神经元}；它们
沿\omterm{def:g8:nervous-communication:neuron}{神经}纤维抵达\emph{\omterm{def:g5:brain-in-command:system}{脊髓}}，在那里直接与运动\omterm{def:g8:nervous-communication:neuron}{神经元}形成\omterm{def:g8:nervous-communication:synapse}{突触}，
随后运动\omterm{def:g8:nervous-communication:neuron}{神经元}的冲动使大腿\omterm{def:g3:bones-and-muscles:muscle}{肌肉}收缩。两类\omterm{def:g8:nervous-communication:neuron}{神经元}、一个\omterm{def:g8:nervous-communication:synapse}{突触}，
均在\omterm{def:g5:brain-in-command:system}{脊髓}水平——
\omterm{def:g5:brain-in-command:system}{大脑}被绕过，因为保护性反应
等不及深思熟虑。

\textbf{2.} 接住物体的回路是从眼睛到脑、再到手臂——
纤维长达数米，并经过许多加权连接；膝跳反射
则是从肌腱到\omterm{def:g5:brain-in-command:system}{脊髓}再到\omterm{def:g3:bones-and-muscles:muscle}{肌肉}——仅厘米量级，
基本只需一个\omterm{def:g8:nervous-communication:synapse}{突触}。线路更短，连接更少，
所需时间仅为十分之一。

\textbf{3.} 决定在\omterm{def:g5:brain-in-command:system}{脊髓}中作出，
早于任何来自脑的否决抵达——且\omterm{def:g8:nervous-communication:synapse}{突触}的单向传递意味着
\omterm{def:g5:brain-in-command:system}{大脑}之后的反对意见无法沿感觉通路
逆向传回。提醒并不能改变什么：
回路在意志之下便已闭合。

\textbf{4.} 感觉是第二段旅程：同一次敲击的
报告也沿\omterm{def:g5:brain-in-command:system}{脊髓}上行至脑——路径更长，
经过更多连接——抵达时，
踢腿早已成为过去。

\textbf{5.} 频率编码：较难的路段发出更快的鼓点，
运动通路以同样的方式回应——
更高的频率带来更强的收缩，
两端同用这一种“单词语言”。

\textbf{6.} 整个回路完好：感觉通路、
\omterm{def:g5:brain-in-command:system}{脊髓}连接、运动通路与\omterm{def:g3:bones-and-muscles:muscle}{肌肉}——双侧对称
也证明两侧同样完好。一次敲击检验了四个
组成部分。

\textbf{7.} 系统其他部分的连接\omterm{def:g8:nervous-communication:neuron}{神经元}
的信号甚至也能到达这一回路的连接处，
提高其准备程度：“无脑”的反射仍然
连在整体之中，其\omterm{def:g8:nervous-communication:synapse}{突触}听到的
比肌腱更多。

\textbf{8.} \omterm{def:g8:nervous-communication:synapse}{突触}——以及纤维的信号传递
总体而言：麻醉剂使连接处的化学过程
平静下来，信息传递也随之减弱，
踢腿与感觉一并消失。

\textbf{9.} 报告（肌腱拉伸）——判定
（\omterm{def:g5:brain-in-command:system}{脊髓}中的单一连接处）——下令
（运动通路）：此回路的中间一步
缩减为单个\omterm{def:g8:nervous-communication:synapse}{突触}，脑的权衡被绕过。

\textbf{10.} 反射：肌腱—\omterm{def:g5:brain-in-command:system}{脊髓}—\omterm{def:g3:bones-and-muscles:muscle}{肌肉}，一个连接处，
五十分之一秒，无法抑制，无需练习。
接住：眼睛—脑—手臂，连接处数不清，
五分之一秒，完全出于自主，需经练习养成。

\textbf{11.} 因为脑的权衡需要时间，而
烧伤每一毫秒都在加深：保护组织的回路
被焊得很短，位于\omterm{def:g5:brain-in-command:system}{脊髓}水平，事后再告知
脑——感觉得晚，此时已获救。

\textbf{12.} “一次锤击展示了这一种单词频率语言、
经\omterm{def:g5:brain-in-command:system}{脊髓}的双线布线，以及一个比其主人
决定得更快的连接处。”
\end{solution}
